% ******************************************************************************
% ****************************** Custom Margin *********************************

% Add `custommargin' in the document class options to use this section
% Set {innerside margin / outerside margin / topmargin / bottom margin}  and
% other page dimensions
\ifsetCustomMargin
  \RequirePackage[left=37mm,right=30mm,top=35mm,bottom=30mm]{geometry}
  \setFancyHdr % To apply fancy header after geometry package is loaded
\fi

% Add spaces between paragraphs
%\setlength{\parskip}{0.5em}
% Ragged bottom avoids extra whitespaces between paragraphs
\raggedbottom
% To remove the excess top spacing for enumeration, list and description
%\usepackage{enumitem}
%\setlist[enumerate,itemize,description]{topsep=0em}

% *****************************************************************************
% ******************* Fonts (like different typewriter fonts etc.)*************

% Add `customfont' in the document class option to use this section

\ifsetCustomFont
  % Set your custom font here and use `customfont' in options. Leave empty to
  % load computer modern font (default LaTeX font).
  %\RequirePackage{helvet}

  % For use with XeLaTeX
  %  \setmainfont[
  %    Path              = ./libertine/opentype/,
  %    Extension         = .otf,
  %    UprightFont = LinLibertine_R,
  %    BoldFont = LinLibertine_RZ, % Linux Libertine O Regular Semibold
  %    ItalicFont = LinLibertine_RI,
  %    BoldItalicFont = LinLibertine_RZI, % Linux Libertine O Regular Semibold Italic
  %  ]
  %  {libertine}
  %  % load font from system font
  %  \newfontfamily\libertinesystemfont{Linux Libertine O}
\fi

% *****************************************************************************
% **************************** Custom Packages ********************************

% ************************* Algorithms and Pseudocode **************************

%\usepackage{algpseudocode}

% *************************** Nomenclature ************************************
% \usepackage{nomencl}

% ********************Captions and Hyperreferencing / URL **********************

% Captions: This makes captions of figures use a boldfaced small font.
%\RequirePackage[small,bf]{caption}

\RequirePackage[labelsep=space,tableposition=top]{caption}
\renewcommand{\figurename}{Fig.} %to support older versions of captions.sty


% *************************** Graphics and figures *****************************

%\usepackage{rotating}
%\usepackage{wrapfig}

% Uncomment the following two lines to force Latex to place the figure.
% Use [H] when including graphics. Note 'H' instead of 'h'
%\usepackage{float}
%\restylefloat{figure}

% Subcaption package is also available in the sty folder you can use that by
% uncommenting the following line
% This is for people stuck with older versions of texlive
%\usepackage{sty/caption/subcaption}
\usepackage{subcaption}

% ********************************** Tables ************************************
\usepackage{booktabs} % For professional looking tables
\usepackage{multirow}

%\usepackage{multicol}
%\usepackage{longtable}
%\usepackage{tabularx}


% *********************************** SI Units *********************************
\usepackage{siunitx} % use this package module for SI units


% ******************************* Line Spacing *********************************

% Choose linespacing as appropriate. Default is one-half line spacing as per the
% University guidelines

% \doublespacing
% \onehalfspacing
% \singlespacing


% ************************ Formatting / Footnote *******************************

% Don't break enumeration (etc.) across pages in an ugly manner (default 10000)
%\clubpenalty=500
%\widowpenalty=500

%\usepackage[perpage]{footmisc} %Range of footnote options


% *****************************************************************************
% *************************** Bibliography  and References ********************

%\usepackage{cleveref} %Referencing without need to explicitly state fig /table

% Add `custombib' in the document class option to use this section
% \ifuseCustomBib
    %    \RequirePackage[square, sort, numbers, authoryear]{natbib} % CustomBib

    % If you would like to use biblatex for your reference management, as opposed to the default `natbibpackage` pass the option `custombib` in the document class. Comment out the previous line to make sure you don't load the natbib package. Uncomment the following lines and specify the location of references.bib file

\RequirePackage[backref=true, style=authoryear, bibstyle=trad-abbrv, citestyle=authoryear, sorting=nty,uniquelist=false,maxcitenames=1]{biblatex}

% Custom citation formatting for author-year style
\renewcommand{\nameyeardelim}{\addcomma\space}
\renewcommand{\multicitedelim}{\addsemicolon\space}
% Force maximum of 1 name in citations
\ExecuteBibliographyOptions{maxcitenames=1, minnames=1}
\DeclareCiteCommand{\cite}[\mkbibparens]
  {\usebibmacro{prenote}}
  {\usebibmacro{citeindex}\usebibmacro{cite}}
  {\multicitedelim}
  {\usebibmacro{postnote}}
% Custom bibliography spacing
\setlength{\bibparsep}{0.5\baselineskip}

\addbibresource{References/references.bib} %Location of references.bib only for biblatex, Do not omit the .bib extension from the filename.
\AtEveryBibitem{\iffieldundef{label}{}{\clearfield{label}}}

% \fi

% changes the default name `Bibliography` -> `References'
\renewcommand{\bibname}{References}


% ******************************************************************************
% ************************* User Defined Commxands ******************************
% ******************************************************************************

% *********** To change the name of Table of Contents / LOF and LOT ************

%\renewcommand{\contentsname}{My Table of Contents}
%\renewcommand{\listfigurename}{My List of Figures}
%\renewcommand{\listtablename}{My List of Tables}


% ********************** TOC depth and numbering depth *************************

\setcounter{secnumdepth}{2}
\setcounter{tocdepth}{2}


% ******************************* Nomenclature *********************************

% To change the name of the Nomenclature section, uncomment the following line

\renewcommand{\nomname}{Symbols}


% ********************************* Appendix ***********************************

% The default value of both \appendixtocname and \appendixpagename is `Appendices'. These names can all be changed via:

%\renewcommand{\appendixtocname}{List of appendices}
%\renewcommand{\appendixname}{Appndx}

% *********************** Configure Draft Mode **********************************

% Uncomment to disable figures in `draft'
%\setkeys{Gin}{draft=true}  % set draft to false to enable figures in `draft'

% These options are active only during the draft mode
% Default text is "Draft"
%\SetDraftText{DRAFT}

% Default Watermark location is top. Location (top/bottom)
%\SetDraftWMPosition{bottom}

% Draft Version - default is v1.0
%\SetDraftVersion{v1.1}

% Draft Text grayscale value (should be between 0-black and 1-white)
% Default value is 0.75
%\SetDraftGrayScale{0.8}


% ******************************** Todo Notes **********************************
%% Uncomment the following lines to have todonotes.

% \ifsetDraft
\usepackage[colorinlistoftodos,prependcaption,textsize=tiny]{todonotes}
\newcommand{\mynote}[1]{\todo[author=sp,size=\small,inline,color=green!40]{#1}}
% \else
% 	\newcommand{\mynote}[1]{}
% 	\newcommand{\listoftodos}{}
% \fi

% Example todo: \mynote{Hey! I have a note}

% ******************************** Highlighting Changes **********************************
%% Uncomment the following lines to be able to highlight text/modifications.
%\ifsetDraft
%  \usepackage{color, soul}
%  \newcommand{\hlc}[2][yellow]{{\sethlcolor{#1} \hl{#2}}}
%  \newcommand{\hlfix}[2]{\texthl{#1}\todo{#2}}
%\else
%  \newcommand{\hlc}[2]{}
%  \newcommand{\hlfix}[2]{}
%\fi

% Example highlight 1: \hlc{Text to be highlighted}
% Example highlight 2: \hlc[green]{Text to be highlighted in green colour}
% Example highlight 3: \hlfix{Original Text}{Fixed Text}

% *****************************************************************************
% ******************* Better enumeration my MB*************
\usepackage{enumitem}
\usepackage{algorithm}
\usepackage{algorithmic}
\usepackage{xltabular}
\usepackage{hyperref}
\usepackage{microtype}
\usepackage{lettrine}
\usepackage{standalone}

\renewcommand{\floatpagefraction}{.9}    % require fuller float pages
\renewcommand{\topfraction}{.9}          % fraction of page allowed for floats at top

\RequirePackage{lmodern}
\usepackage[commands,enumerate,environments,notes]{SP}

\newcommand{\Dtrain}{\mathcal{D}_{\text{train}}}

% \newcommand{\E}{\mathbb{E}}

\newcommand{\dd}{\mathrm{d}}

\newcommand{\hpi}{\hat{\pi}}

\newcommand{\piy}{\pi_{\cdot \mid y}}

\newcommand{\Xt}{\{X_t\}_{t \in T}} 

\newcommand{\rightP}{\overrightarrow{\mathbb{P}}}
\newcommand{\leftP}{\overleftarrow{\mathbb{P}}}


% --- linea for captions ----------------------------------------------
\usepackage{tikz,xcolor}

% helper that keeps the line vertically centred on the baseline
\newcommand{\RaiseLine}[1]{\raisebox{.5ex}{\tikz{#1}}}

% ▬ ▬ thick black, widely‑spaced dashes (first row in your figure)
\newcommand{\thickblackdash}{%
  \RaiseLine{%
    \draw[black,line width = 1.6pt,dash pattern = on 5pt off 3pt,line cap=butt] (0,0) -- (5mm,0);}}

% – – – thin grey, evenly‑spaced dashes (second row in your figure)
\newcommand{\thingreydash}{%
  \RaiseLine{%
    \draw[gray!60,line width = 0.9pt,dash pattern = on 3pt off 3pt,line cap=round] (0,0) -- (5mm,0);}}
% -----------------------------------------------------------------------------


% ---------- a tiny helper ----------
% \badge{<bg-colour>}{<stuff>}  →  coloured box wrapping <stuff> in display-style math
\newcommand{\badge}[2]{\colorbox{#1}{\ensuremath{#2}}}

\newcommand{\Ra}{\mathcal{R}a}
\newcommand{\Rb}{\mathcal{R}b}

% Macros for forward and backward processes with superscripts
\newcommand{\fwdP}[2]{\overrightarrow{\mathbb{P}}^{#1, #2}}
\newcommand{\bwdP}[2]{\overleftarrow{\mathbb{P}}^{#1, #2}}

% \usepackage{pdflscape}      % provides the landscape environment
% \usepackage{booktabs}       % better horizontal rules
% \usepackage{adjustbox}      % easy table-resizing
% \usepackage{rotating}